% ECONOMICS_PROBLEM: {"id": "F13-565", "title": "Trade-policy uncertainty", "jel_code": "F13", "source_id": "OP-0085"}
% !TeX program = xelatex
\documentclass[11pt,a4paper]{article}
\usepackage[margin=23mm]{geometry}
\usepackage{fontspec}
\setmainfont{Latin Modern Roman}
\usepackage{amsmath,amssymb,mathtools}
\usepackage{xcolor,enumitem,booktabs,longtable,array,xurl,hyperref}
\definecolor{muted}{HTML}{53636C}
\hypersetup{colorlinks=true,urlcolor=blue,linkcolor=blue}
\setlength{\parindent}{0pt}
\setlength{\parskip}{5pt}
\newcommand{\R}{\mathbb R}
\newcommand{\E}{\mathbb E}
\newcommand{\N}{\mathbb N}
\newcommand{\ind}{\mathbf 1}
\newcommand{\Var}{\operatorname{Var}}
\newcommand{\Cov}{\operatorname{Cov}}
\newcommand{\supp}{\operatorname{supp}}
\DeclareMathOperator*{\argmin}{arg\,min}
\DeclareMathOperator*{\argmax}{arg\,max}
\newcommand{\entrymeta}[3]{{\sffamily\small\color{muted}\textbf{\texttt{#1}}\quad|\quad #2\par #3\par}}
\newcommand{\Problemref}[1]{\texttt{#1}}
\newcommand{\JELref}[2]{\texttt{#2} (JEL \texttt{#1})}
\begin{document}
\section*{Trade-policy uncertainty}
\textbf{Problem ID:} \texttt{F13-565}\quad \textbf{JEL:} \texttt{F13}\par
% BEGIN ECONOMICS STATEMENT
\label{problem:OP-0085}
\label{atlas:prob:T5}

\noindent\textbf{Atlas ID:} T5; \textbf{original number:} 85; \textbf{field:} International trade. \textbf{Classification:} Editorial research formulation (theoretical/empirical); not a certified unsolved theorem.

\subsection*{Definitions and assumptions}
At dates $0,\ldots,H$, a firm chooses investment, export entry and sourcing under sunk costs, borrowing limits and tariff path $T=(T_1,\ldots,T_H)$. Model $\theta$ specifies technology, demand, discounting and information. At date zero the firm’s posterior law $\mu$ over future tariffs is part of the state. Other exogenous shocks have a specified joint law with $T$. Define investment $I_\theta(\mu)$ through optimal feasible dynamic choices, jointly with a stated market equilibrium and selection where relevant. Assume finite expected profits and a measurable optimal policy.

\subsection*{Formal research question}
For probability laws $\mu_0,\mu_1$ with common date-by-date expected tariffs and common current tariffs, define
\[
 U_I(\theta;\mu_1,\mu_0)
 =\mathbb E_{\mu_1}[I_\theta(\mu_1)]
  -\mathbb E_{\mu_0}[I_\theta(\mu_0)].
\]
Specify the expectation over remaining observable firm states. Require $\mu_1$ to be a mean-preserving spread of $\mu_0$ in a stated multivariate convex order if that is the intended uncertainty increase. Characterize when $U_I$ is negative, zero or positive, and identify its magnitude separately from changes in expected tariffs, demand, financing conditions and risk premia. Establish what belief surveys, policy announcements and investment panels identify rather than simply correlate with this effect.

\subsection*{Interpretation and scope}
A variance increase is not a complete ordering of multi-period uncertainty, and a credible agreement usually changes both the mean and the distribution of tariffs. Irreversibility can generate waiting incentives, but a universal negative response should not be asserted without restrictions. Distinguish subjective firm beliefs from a researcher’s uncertainty index. Empirical identification needs exogenous changes in the relevant posterior law or an explicit structural restriction connecting announcements, beliefs and choices. For general equilibrium effects, allow entry and input-price feedback rather than summing partial-equilibrium firm responses. Any instrument must affect beliefs through the modeled announcement while excluding unmodeled direct effects on investment demand.

\subsection*{Source audit and status}
[S1] develops policy-uncertainty and sunk-cost mechanisms; [S2] relates a historical change in tariff risk to manufacturing; [S3] constructs uncertainty measures and aggregate evidence. They do not separately identify every mean-preserving uncertainty counterfactual. The question is an editorial causal and comparative-statics agenda whose current completeness is not certified by these references.

\subsection*{References}

\begin{enumerate}[label={[S\arabic*]},leftmargin=*]

\item Kyle Handley and Nuno Limão (2015). \emph{Trade and Investment under Policy Uncertainty: Theory and Firm Evidence}. American Economic Journal: Economic Policy 7(4), 189–222. \url{https://doi.org/10.1257/pol.20140068}.
\textit{Locator and role:} Publisher or author abstract / bibliographic record; Sunk export costs, policy uncertainty and Portugal’s EC accession. Provides background or a benchmark result; does not establish that the editorial question remains unresolved in 2026.

\item Justin R. Pierce and Peter K. Schott (2016). \emph{The Surprisingly Swift Decline of US Manufacturing Employment}. American Economic Review 106(7), 1632–1662. \url{https://doi.org/10.1257/aer.20131578}.
\textit{Locator and role:} Publisher or author abstract / bibliographic record; Changes in tariff risk and US manufacturing employment; not a mean-preserving uncertainty experiment. Provides background or a benchmark result; does not establish that the editorial question remains unresolved in 2026.

\item Dario Caldara, Matteo Iacoviello, Patrick Molligo, Andrea Prestipino, and Andrea Raffo (2020). \emph{The Economic Effects of Trade Policy Uncertainty}. Journal of Monetary Economics 109, 38–59. \url{https://doi.org/10.1016/j.jmoneco.2019.11.002}.
\textit{Locator and role:} Publisher or author abstract / bibliographic record; Trade-policy uncertainty measures and aggregate investment responses. Provides background or a benchmark result; does not establish that the editorial question remains unresolved in 2026.

\end{enumerate}

\subsection*{Common conventions: Source mathematical conventions}

Unless an entry states otherwise, agent and firm sets are finite. State and action spaces are measurable, strategies and outcome maps are measurable, probability laws are specified, and expectations exist for the targets under discussion. Infinite-horizon discount factors lie in $(0,1)$ and discounted objectives are finite. Norms, tolerances and complexity input sizes must be specified for computational comparisons. Constants or primitives that appear locally are inputs, not empirical facts asserted by this catalogue.

\textbf{Identification.} A statistical model $m\in\mathcal M$ implies an observable law $P_m$ and target $\theta(m)$. At law $P$, its identified set is
\[
\Theta(P;\mathcal M)=\{\theta(m):m\in\mathcal M,\ P_m=P\}.
\]
Point identification holds when this set is a singleton, or equivalently when $P_{m_1}=P_{m_2}$ implies $\theta(m_1)=\theta(m_2)$ for all compatible models. Sharp bounds mean bounds attained, or approached when the boundary is not attained, by compatible models. Writing this set is a definition, not a solution: the substantive task is to establish informative restrictions and derive or estimate the set. An unrestricted-confounding model may give only trivial bounds.

\textbf{Inference.} Identification, estimability and valid uncertainty quantification are separate questions. Fix $\alpha\in(0,1)$. A confidence procedure $C_n$ over a declared law class $\mathcal P$ has asymptotic uniform coverage at level $1-\alpha$ when
\[
\liminf_{n\to\infty}\inf_{P\in\mathcal P}\Pr_P\{\theta(P)\in C_n\}\geq1-\alpha.
\]
For set-identified targets the coverage event must be defined separately, for example coverage of the full identified set. If an entry uses identification-indexed models instead of uniquely indexed laws, the same coverage convention is applied over those models. Pointwise limits do not establish uniform validity, and asymptotic validity does not guarantee finite-sample precision. Sample sizes, dependence structures and weak-identification sequences are part of each application.

\textbf{Counterfactuals.} $Y_i(a)$ is a potential outcome under a specified intervention, including its timing, implementation and versions. It is not $Y_i$ conditional on voluntarily choosing $a$. A full intervention vector permits interference; shorthand scalar treatment notation does not silently impose no interference. Assignment, selection, attrition, anticipation and measurement must be specified. A claim about an instrument's local effect is not automatically a population policy effect.

\textbf{Economic equilibrium.} A counterfactual model includes best responses, constraints, feasibility, market clearing, state transitions and expectations consistent with the stated information. When several equilibria exist, either a declared selector is an input or the target is set-valued. Comparisons across policy regimes must use the stated selection convention. Reduced-form relationships are not presumed invariant after an equilibrium-changing intervention.

\textbf{Optimization and welfare.} Optimization is over the stated feasible set. If attainment is not established, $\sup$ or $\inf$ and $\varepsilon$-optimal choices replace an asserted maximizer or minimizer. For example, with value $V=\sup_{a\in\mathcal A}W(a)<\infty$, an $\varepsilon$-optimal set is $\{a\in\mathcal A:W(a)\geq V-\varepsilon\}$ for $\varepsilon>0$. Benefits, costs, risk and health or environmental outcomes can be aggregated only using declared units and conversion weights. Interpersonal utility weights, the treatment of fiscal transfers and foreign profits, discounting, and the behavioral welfare criterion are explicit inputs. A comparative-static derivative additionally requires existence and appropriate differentiability of the selected counterfactual.

\textbf{Sources.} Local labels [S1], [S2], and so forth restart in every question. Locators identify the relevant abstract, model, section or result at the level actually checked; a bibliographic or abstract-level verification is not represented as inspection of an entire proof. Working papers are distinguished from later publications and changing versions. Source summaries are paraphrased. All URL checks and editorial formulations refer to the audit date stated above.
% END ECONOMICS STATEMENT
\end{document}
